\documentclass[10pt,a4paper]{amsart}
\usepackage{amsmath,amssymb,amsthm,graphicx}
\usepackage{verbatim}
\newenvironment{reference}{\comment}{\endcomment}
%\newenvironment{reference}{}{}
\newenvironment{slogan}{\comment}{\endcomment}
\newenvironment{history}{\comment}{\endcomment}
% For commutative diagrams we use Xy-pic
\usepackage[all]{xy}

% We use 2cell for 2-commutative diagrams.
\xyoption{2cell}
% Theorem environments.
%
\theoremstyle{plain}
\newtheorem{theorem}[subsection]{Theorem}
\newtheorem{proposition}[subsection]{Proposition}
\newtheorem{lemma}[subsection]{Lemma}

\theoremstyle{definition}
\newtheorem{definition}[subsection]{Definition}
\newtheorem{example}[subsection]{Example}
\newtheorem{exercise}[subsection]{Exercise}
\newtheorem{situation}[subsection]{Situation}

\theoremstyle{remark}
\newtheorem{remark}[subsection]{Remark}
\newtheorem{remarks}[subsection]{Remarks}

\numberwithin{equation}{subsection}

% Macros
%
\def\lim{\mathop{\mathrm{lim}}\nolimits}
\def\colim{\mathop{\mathrm{colim}}\nolimits}
\def\Spec{\mathop{\mathrm{Spec}}}
\def\Hom{\mathop{\mathrm{Hom}}\nolimits}
\def\Ext{\mathop{\mathrm{Ext}}\nolimits}
\def\SheafHom{\mathop{\mathcal{H}\!\mathit{om}}\nolimits}
\def\SheafExt{\mathop{\mathcal{E}\!\mathit{xt}}\nolimits}
\def\Sch{\mathit{Sch}}
\def\Mor{\mathop{\mathrm{Mor}}\nolimits}
\def\Ob{\mathop{\mathrm{Ob}}\nolimits}
\def\Sh{\mathop{\mathit{Sh}}\nolimits}
\def\NL{\mathop{N\!L}\nolimits}
\def\CH{\mathop{\mathrm{CH}}\nolimits}
\def\proetale{{pro\text{-}\acute{e}tale}}
\def\etale{{\acute{e}tale}}
\def\QCoh{\mathit{QCoh}}
\def\Ker{\mathop{\mathrm{Ker}}}
\def\Im{\mathop{\mathrm{Im}}}
\def\Coker{\mathop{\mathrm{Coker}}}
\def\Coim{\mathop{\mathrm{Coim}}}

% Boxtimes
%
\DeclareMathSymbol{\boxtimes}{\mathbin}{AMSa}{"02}

%
% Macros for moduli stacks/spaces
%
\def\QCohstack{\mathcal{QC}\!\mathit{oh}}
\def\Cohstack{\mathcal{C}\!\mathit{oh}}
\def\Spacesstack{\mathcal{S}\!\mathit{paces}}
\def\Quotfunctor{\mathrm{Quot}}
\def\Hilbfunctor{\mathrm{Hilb}}
\def\Curvesstack{\mathcal{C}\!\mathit{urves}}
\def\Polarizedstack{\mathcal{P}\!\mathit{olarized}}
\def\Complexesstack{\mathcal{C}\!\mathit{omplexes}}
% \Pic is the operator that assigns to X its picard group, usage \Pic(X)
% \Picardstack_{X/B} denotes the Picard stack of X over B
% \Picardfunctor_{X/B} denotes the Picard functor of X over B
\def\Pic{\mathop{\mathrm{Pic}}\nolimits}
\def\Picardstack{\mathcal{P}\!\mathit{ic}}
\def\Picardfunctor{\mathrm{Pic}}
\def\Deformationcategory{\mathcal{D}\!\mathit{ef}}
\usepackage{cleveref}
\crefname{originalproposition}{original proposition}{original propositions}
\Crefname{originalproposition}{Original proposition}{Original propositions}
\makeatletter
\newlabel{cohomology-proposition-cohomological-dimension-spectral}{{0A3G}{source}{}{sourcepointer-1}{}}
\newlabel{cohomology-proposition-cohomological-dimension-spectral@cref}{{[originalproposition][0][]0A3G}{[1][0][]source}{}{}{}}
\makeatother
\crefname{originallemma}{original lemma}{original lemmas}
\Crefname{originallemma}{Original lemma}{Original lemmas}
\makeatletter
\newlabel{decent-spaces-lemma-dimension-decent-space}{{0A4J}{source}{}{sourcepointer-2}{}}
\newlabel{decent-spaces-lemma-dimension-decent-space@cref}{{[originallemma][0][]0A4J}{[1][0][]source}{}{}{}}
\makeatother
\crefname{originallemma}{original lemma}{original lemmas}
\Crefname{originallemma}{Original lemma}{Original lemmas}
\makeatletter
\newlabel{decent-spaces-lemma-filter-quasi-compact-quasi-separated}{{07ST}{source}{}{sourcepointer-3}{}}
\newlabel{decent-spaces-lemma-filter-quasi-compact-quasi-separated@cref}{{[originallemma][0][]07ST}{[1][0][]source}{}{}{}}
\makeatother
\crefname{originalproposition}{original proposition}{original propositions}
\Crefname{originalproposition}{Original proposition}{Original propositions}
\makeatletter
\newlabel{descent-proposition-same-cohomology-quasi-coherent}{{03DW}{source}{}{sourcepointer-4}{}}
\newlabel{descent-proposition-same-cohomology-quasi-coherent@cref}{{[originalproposition][0][]03DW}{[1][0][]source}{}{}{}}
\makeatother
\crefname{originallemma}{original lemma}{original lemmas}
\Crefname{originallemma}{Original lemma}{Original lemmas}
\makeatletter
\newlabel{etale-cohomology-lemma-cohomology-with-support-quasi-coherent}{{0A46}{source}{}{sourcepointer-5}{}}
\newlabel{etale-cohomology-lemma-cohomology-with-support-quasi-coherent@cref}{{[originallemma][0][]0A46}{[1][0][]source}{}{}{}}
\makeatother
\crefname{originallemma}{original lemma}{original lemmas}
\Crefname{originallemma}{Original lemma}{Original lemmas}
\makeatletter
\newlabel{lemma-etale-localization-sheaf-with-support}{{0A4P}{source}{}{sourcepointer-6}{}}
\newlabel{lemma-etale-localization-sheaf-with-support@cref}{{[originallemma][0][]0A4P}{[1][0][]source}{}{}{}}
\makeatother
\crefname{originallemma}{original lemma}{original lemmas}
\Crefname{originallemma}{Original lemma}{Original lemmas}
\makeatletter
\newlabel{spaces-properties-lemma-dimension-decent-invariant-under-etale}{{0A4H}{source}{}{sourcepointer-7}{}}
\newlabel{spaces-properties-lemma-dimension-decent-invariant-under-etale@cref}{{[originallemma][0][]0A4H}{[1][0][]source}{}{}{}}
\makeatother
\title{Vanishing of cohomology above the dimension of a quasi-compact quasi-separated algebraic space, with surjectivity in the top degree and vanishing with support}
\author{The Stacks Project Authors}
\date{stacks-project:spaces-cohomology.tex:local-snapshot-20261012; GNU FDL 1.2 or later}
\begin{document}
\maketitle
\noindent\textit{Editorial note.} All delivered mathematical text is the author's own, copied from the pinned Stacks Project chapter file spaces-cohomology.tex; no statement or proof step was rewritten, reordered or removed.\par\medskip
\noindent\textit{Editorial note.} Preamble provenance: the pinned chapter file opens with \textbackslash{}input\{preamble\} and carries no preamble text of its own, so the project's single Stacks-wide preamble staged at sources/stacks-spaces-cohomology/preamble.tex is excerpted verbatim as three spans, lines 8-12 (verbatim package and the project's own reference/slogan/history display environments), lines 14-18 (Xy-pic setup) and lines 154-217 (theorem declarations and mathematical macros). The wrapper class and the amsmath/amssymb/amsthm/graphicx base packages are supplied by the assembler. No macro was invented or redefined.\par\medskip
\noindent\textit{Editorial note.} Same-chapter pointers are located at their \textbackslash{}label lines inside the pinned chapter file; cross-chapter pointers are located in the project's own tag list staged in the snapshot (sources/stacks-chow/tags.txt), at the verbatim lines carrying their canonical tags.\par\medskip
\noindent\textit{Allowed standard prerequisites (editorial).}\begin{itemize}
\item Dimension is invariant under etale morphisms, and for a quasi-separated decent space it equals the Krull dimension. Properties of Spaces, label spaces-properties-lemma-dimension-decent-invariant-under-etale (invoked at proof line 1785) and Decent Spaces, label decent-spaces-lemma-dimension-decent-space (invoked at proof line 1789). They let the argument work with honest Krull dimension on the auxiliary schemes that appear, everything etale over \$X\$ having dimension at most \$d\$.
\item Etale cohomology of a quasi-coherent module on a scheme agrees with usual cohomology, also with support in a closed subscheme. Descent, label descent-proposition-same-cohomology-quasi-coherent (invoked at proof line 1795) and Etale Cohomology, label etale-cohomology-lemma-cohomology-with-support-quasi-coherent (invoked at proof line 1798). They identify the cohomology computed on the auxiliary schemes \$V\_p\$ with the usual sheaf cohomology, where the vanishing statement for schemes applies.
\item Filtration of a quasi-compact quasi-separated algebraic space by open pieces with etale schemes over them. Decent Spaces, label decent-spaces-lemma-filter-quasi-compact-quasi-separated (invoked at proof line 1803). It provides the chain \$\textbackslash{}emptyset = U\_\{n+1\} \textbackslash{}subset U\_n \textbackslash{}subset \textbackslash{}ldots \textbackslash{}subset U\_1 = X\$ and the quasi-compact separated schemes \$V\_p \textbackslash{}to U\_p\$ which are surjective etale and isomorphisms over \$T\_p = U\_p \textbackslash{}setminus U\_\{p+1\}\$.
\item Vanishing of cohomology above the cohomological dimension of a scheme. Cohomology, label cohomology-proposition-cohomological-dimension-spectral (invoked at proof lines 1818 and 1831). It gives the vanishing over the auxiliary schemes \$V\_p\$, which is transported back along the filtration pieces.
\item Localization of cohomology with support along an etale morphism that is an isomorphism over the support. Same chapter spaces-cohomology.tex, label lemma-etale-localization-sheaf-with-support (invoked at proof lines 1828 and 1839). It converts the cohomology with support \$H\textasciicircum{}q\_\{T\_p\}(U\_p, \textbackslash{}mathcal\{F\})\$ into cohomology with support on the scheme \$V\_p\$, where vanishing is known.
\end{itemize}\medskip
\ifdefined\hypertarget\hypertarget{sourcepointer-1}{}\fi
\noindent\textit{Source reference.} Original proposition 0A3G: sources/stacks-chow/tags.txt, source lines 12739--12739. Navigation pointer for \textbackslash{}ref\{cohomology-proposition-cohomological-dimension-spectral\} used inside the counted span (line 1818); the argument that consumes it is retained in the span.\par\medskip
\ifdefined\hypertarget\hypertarget{sourcepointer-2}{}\fi
\noindent\textit{Source reference.} Original lemma 0A4J: sources/stacks-chow/tags.txt, source lines 12777--12777. Navigation pointer for \textbackslash{}ref\{decent-spaces-lemma-dimension-decent-space\} used inside the counted span (line 1789); the argument that consumes it is retained in the span.\par\medskip
\ifdefined\hypertarget\hypertarget{sourcepointer-3}{}\fi
\noindent\textit{Source reference.} Original lemma 07ST: sources/stacks-chow/tags.txt, source lines 9913--9913. Navigation pointer for \textbackslash{}ref\{decent-spaces-lemma-filter-quasi-compact-quasi-separated\} used inside the counted span (line 1803); the argument that consumes it is retained in the span.\par\medskip
\ifdefined\hypertarget\hypertarget{sourcepointer-4}{}\fi
\noindent\textit{Source reference.} Original proposition 03DW: sources/stacks-chow/tags.txt, source lines 4435--4435. Navigation pointer for \textbackslash{}ref\{descent-proposition-same-cohomology-quasi-coherent\} used inside the counted span (line 1795); the argument that consumes it is retained in the span.\par\medskip
\ifdefined\hypertarget\hypertarget{sourcepointer-5}{}\fi
\noindent\textit{Source reference.} Original lemma 0A46: sources/stacks-chow/tags.txt, source lines 12764--12764. Navigation pointer for \textbackslash{}ref\{etale-cohomology-lemma-cohomology-with-support-quasi-coherent\} used inside the counted span (line 1798); the argument that consumes it is retained in the span.\par\medskip
\ifdefined\hypertarget\hypertarget{sourcepointer-6}{}\fi
\noindent\textit{Source reference.} Original lemma 0A4P: sources/stacks-spaces-cohomology/pinned-original-source.tex, source lines 1674--1674. Navigation pointer for \textbackslash{}ref\{lemma-etale-localization-sheaf-with-support\} used inside the counted span (line 1828); the argument that consumes it is retained in the span.\par\medskip
\ifdefined\hypertarget\hypertarget{sourcepointer-7}{}\fi
\noindent\textit{Source reference.} Original lemma 0A4H: sources/stacks-chow/tags.txt, source lines 12775--12775. Navigation pointer for \textbackslash{}ref\{spaces-properties-lemma-dimension-decent-invariant-under-etale\} used inside the counted span (line 1785); the argument that consumes it is retained in the span.\par\medskip
\begin{lemma}
\label{lemma-vanishing-above-dimension}
Let $S$ be a scheme. Let $X$ be a quasi-compact and quasi-separated
algebraic space over $S$. Assume $\dim(X) \leq d$ for some integer $d$.
Let $\mathcal{F}$ be a quasi-coherent sheaf $\mathcal{F}$ on $X$.
\begin{enumerate}
\item $H^q(X, \mathcal{F}) = 0$ for $q > d$,
\item $H^d(X, \mathcal{F}) \to H^d(U, \mathcal{F})$ is surjective
for any quasi-compact open $U \subset X$,
\item $H^q_Z(X, \mathcal{F}) = 0$ for $q > d$ for any closed subspace
$Z \subset X$ whose complement is quasi-compact.
\end{enumerate}
\end{lemma}

\begin{proof}
By Properties of Spaces, Lemma
\ref{spaces-properties-lemma-dimension-decent-invariant-under-etale}
every algebraic space $Y$ \'etale over $X$ has dimension $\leq d$.
If $Y$ is quasi-separated, the dimension of $Y$ is equal to the
Krull dimension of $|Y|$ by 
Decent Spaces, Lemma \ref{decent-spaces-lemma-dimension-decent-space}.
Also, if $Y$ is a scheme, then \'etale cohomology of $\mathcal{F}$
over $Y$, resp.\ \'etale cohomology of $\mathcal{F}$ with support in a
closed subscheme, agrees with usual cohomology of $\mathcal{F}$,
resp.\ usual cohomology with support in the closed subscheme.
See
Descent, Proposition \ref{descent-proposition-same-cohomology-quasi-coherent}
and
\'Etale Cohomology, Lemma
\ref{etale-cohomology-lemma-cohomology-with-support-quasi-coherent}.
We will use these facts without further mention.

\medskip\noindent
By Decent Spaces, Lemma
\ref{decent-spaces-lemma-filter-quasi-compact-quasi-separated}
there exist an integer $n$ and open subspaces
$$
\emptyset = U_{n + 1} \subset
U_n \subset U_{n - 1} \subset \ldots \subset U_1 = X
$$
with the following property: setting $T_p = U_p \setminus U_{p + 1}$
(with reduced induced subspace structure) there exists a quasi-compact
separated scheme $V_p$ and a surjective \'etale morphism $f_p : V_p \to U_p$
such that $f_p^{-1}(T_p) \to T_p$ is an isomorphism.

\medskip\noindent
As $U_n = V_n$ is a scheme, our initial remarks imply the cohomology of
$\mathcal{F}$ over $U_n$ vanishes in degrees $> d$ by
Cohomology, Proposition
\ref{cohomology-proposition-cohomological-dimension-spectral}.
Suppose we have shown, by induction, that
$H^q(U_{p + 1}, \mathcal{F}|_{U_{p + 1}}) = 0$ for $q > d$.
It suffices to show $H_{T_p}^q(U_p, \mathcal{F})$ for
$q > d$ is zero in order to conclude the vanishing of cohomology
of $\mathcal{F}$ over $U_p$ in degrees $> d$.
However, we have
$$
H^q_{T_p}(U_p, \mathcal{F}) = H^q_{f_p^{-1}(T_p)}(V_p, \mathcal{F})
$$
by Lemma \ref{lemma-etale-localization-sheaf-with-support}
and as $V_p$ is a scheme we obtain the desired vanishing from
Cohomology, Proposition
\ref{cohomology-proposition-cohomological-dimension-spectral}.
In this way we conclude that (1) is true.

\medskip\noindent
To prove (2) let $U \subset X$ be a quasi-compact open subspace.
Consider the open subspace $U' = U \cup U_n$. Let $Z = U' \setminus U$.
Then $g : U_n \to U'$ is an \'etale morphism such that
$g^{-1}(Z) \to Z$ is an isomorphism. Hence by
Lemma \ref{lemma-etale-localization-sheaf-with-support}
we have $H^q_Z(U', \mathcal{F}) = H^q_Z(U_n, \mathcal{F})$
which vanishes in degree $> d$ because $U_n$ is a scheme
and we can apply
Cohomology, Proposition
\ref{cohomology-proposition-cohomological-dimension-spectral}.
We conclude that $H^d(U', \mathcal{F}) \to H^d(U, \mathcal{F})$
is surjective. Assume, by induction, that we have reduced
our problem to the case where $U$ contains $U_{p + 1}$.
Then we set $U' = U \cup U_p$, set $Z = U' \setminus U$, and
we argue using the morphism $f_p : V_p \to U'$ which is \'etale
and has the property that $f_p^{-1}(Z) \to Z$ is an isomorphism.
In other words, we again see that
$$
H^q_Z(U', \mathcal{F}) = H^q_{f_p^{-1}(Z)}(V_p, \mathcal{F})
$$
and we again see this vanishes in degrees $> d$.
We conclude that $H^d(U', \mathcal{F}) \to H^d(U, \mathcal{F})$
is surjective. Eventually we reach the stage where $U_1 = X \subset U$
which finishes the proof.

\medskip\noindent
A formal argument shows that (2) implies (3).
\end{proof}

\end{document}
